\section{What an effective inquiry can learn}
\label{sec:inquiry}
An effective adaptive protocol has finitely encoded histories, a
computable choice of its next finite experiment and rational output,
and uniformly computable finite conditional outcome laws on valid
positive-probability histories. Every successful stage finishes in
finite time; no computable cost or convergence bound is assumed.
Finite systems with algebraic Hamiltonians and measurements, and
evolution for computable finite times, satisfy these hypotheses:
matrix exponentiation is computable, although its entries need not
remain algebraic. The controller may use independent fair coins.
Noncomputable control parameters, oracle states and an infinite
target-advice tape are outside this model. The conclusion is
conditional on this effective experimental description.

\begin{theorem}[Fixed-target randomized reach]\label{thm:reach}
If such a protocol's rational outputs converge to a fixed real
$\alpha$ with positive probability, then $\alpha\in\mathsf D_2$.
Conversely every $\mathsf D_2$ real is the almost-sure limit of
some such protocol. Consequently a fixed exponent outside
$\mathsf D_2$ in Theorem~\ref{thm:register-headline} has zero
convergence probability for every effective protocol in this class.
\end{theorem}
\begin{proof}
Couple all outcomes to a single fair binary tape $X$. Recursively
partition a unit interval into subintervals in the conditional
probabilities prescribed after each history. Every finite-history
endpoint is a computable real. Outside their countable null set,
strict comparisons are decided after finitely many tape bits and
probability approximations. Interval lengths reproduce all finite
joint laws, including adaptive settings. Direct controller coins
can be included as further steps. Thus outputs $q_k$ are partial
$X$-computable functions defined at every successful stage outside
that null set.

On the measurable event $q_k\to\alpha$, the jump $X'$ computes a
Cauchy name of $\alpha$ by deciding whether a proposed rational tail
contains a deviation beyond a requested tolerance. Almost every
tape is 2-random and satisfies $X'\leq_T X\oplus0'$.
This generalized-lowness theorem is due to Sacks and Stillwell;
see \cite[Corollary~3.12]{randomDegrees} and its attribution through
Kautz's thesis. Intersecting with the success event preserves its
positive measure. If $\alpha$ is irrational its rational lower cut
is computable from its Cauchy name; if it is rational it is already
computable.

We use the following relativized cone argument. If a fixed set $A$
is computable from $X\oplus Z$ on a positive-measure set of tapes,
some fixed functional does so on a positive-measure subset. The
Lebesgue density theorem supplies a finite cylinder on which this
subset has relative measure greater than $1/2$. For each input,
enumerate relative to $Z$ the halting cylinders yielding answer
zero and those yielding answer one. Their measures are lower
semicomputable and disjoint; only the correct answer reaches a
certified measure greater than $1/2$ in the chosen cylinder. Waiting
for it computes $A$ from $Z$. The cylinder and functional are finite
nonuniform constants. This is the usual positive-measure cone
mechanism; see also \cite{probabilisticChoice}. Apply it with
$Z=0'$ and the cut of $\alpha$. Thus $\alpha\leq_T0'$.

Conversely choose a computable rational approximation to a given
$\mathsf D_2$ real. A protocol that prints it and ignores any dummy
experiment succeeds on every run. Finally, a positive success
probability for an exponent outside $\mathsf D_2$ would contradict
the first conclusion.
\end{proof}
The deterministic converse is the classical limit-computability
phenomenon of Gold and Putnam \cite{goldLimit,putnamTrial}.
The randomized obstruction combines generalized lowness for almost
every oracle with the relativized positive-measure cone theorem.
Our use here is their fixed-target formulation for effective adaptive
experiments and its application to the constructed gap exponents.
This target-by-target existential statement is distinct from
identification of an unknown system by one learner, as studied in
formal learning theory \cite{kellyLearning}.

The target is fixed independently of the outcome tape. A computable
protocol can output successive binary prefixes of a random real,
whose realized limit is almost surely outside $\mathsf D_2$; each
fixed target then has probability zero. Likewise fixing an arbitrary
noncomputable tape gives an oracle computation, not an ordinary
deterministic algorithm. The converse above is existential and may
ignore the data; it is not a common learner identifying unknown
systems. Every nonnegative finite $\mathsf D_2$ target also has a
planar material realization here, but no spectroscopy assumption is
needed for that existential statement. Only the exponents
\emph{outside} $\mathsf D_2$, not all of $\mathsf D_3$, have the
stated zero-reach obstruction. Countably many effective protocols
still give a null union of successes for any one such fixed target.

\begin{theorem}[Hidden finite-size scale]\label{thm:hidden-scale}
Suppose finite gaps $\Delta(L,g)$ are uniformly computable on
permitted sizes and rational $0<g<g_0$. Suppose
$D(g)=\lim_L\Delta(L,g)>0$ and the existing rational-coupling
exponent $\alpha=\lim_{g\downarrow0}\log(1/D(g))/\log(1/g)$
is finite. Assume permitted sizes have a computable cofinal
rounding-up selector $s(n)\geq n$. Let $L_*(g)$ be the least
permitted size such that every later permitted size satisfies
\[
 D(g)/2\leq\Delta(L,g)\leq2D(g).
\]
If $L_*(g)$ has a total computable eventual upper bound on rational
couplings near zero, then $\alpha\in\mathsf D_2$.
\end{theorem}
\begin{proof}
Hard-code rational $0<r<\min(g_0,1)$ inside the eventual domain
and take $g_k=r2^{-k-1}$. Given a bound $B$, choose
$L_k=s(\max(1,\lceil B(g_k)\rceil))$. The factor-two inequality
gives normalized-logarithm error at most
$\log2/\log(1/g_k)\to0$. Compute the finite normalized gap to
rational error $2^{-k}$; promised positivity permits a terminating
search for a positive lower enclosure before evaluating its logarithm.
These computable rationals converge to $\alpha$, proving the claim.
For factorial sides compute successive factorials until the requested
threshold is reached. Locality is not used in this implication.
\end{proof}
For a beyond-$\mathsf D_2$ exponent the contrapositive is precise:
for every total computable $B$ and rational $\eta>0$ there is a
rational $0<g<\min(\eta,g_0)$ with $L_*(g)>B(g)$.
More strongly, along any computable rational $g_k\to0$, every
computable bound $b(k)$ has infinitely many violations; otherwise
the same proof would compute a convergent rational approximation.
This is absence of an eventual computable upper bound, not eventual
domination at every coupling or an algorithm locating violations.
The least relative-error tail threshold is also not automatically an
optimal experimental size or a certified stopping rule.
The construction's structural pitch is computable and proportional
to $1/g$; that fact neither bounds the settling threshold nor
identifies the pitch with a correlation length.

\begin{corollary}[Incompatibility with a uniform single-length tail]
\label{cor:uniform-scaling}
The beyond-$\mathsf D_2$ realizations admit no factor-two
thermodynamic tail bound valid at all permitted
$L\geq Cg^{-N}$ for all small rational $g$, with finite constants
$C,N$. The same holds with any computable raw threshold, including
a fixed finite-parameter exponential threshold after enlarging its
constants to rational or integer bounds.
\end{corollary}
\begin{proof}
Round that raw threshold up by $s$. It is then a computable upper
bound for the least \emph{permitted} $L_*$, contrary to
Theorem~\ref{thm:hidden-scale}. For example a raw threshold 25
rounds to factorial side 120, so the rounding must not be omitted.
\end{proof}
In particular consider the explicit ansatz
\[
 \Delta(L,g)=L^{-z}\Phi(L/\xi(g))(1+r(L,g)),\qquad
 \frac{\Phi(x)}{\Phi_\infty x^z}\longrightarrow1,
 \quad D(g)=\Phi_\infty\xi(g)^{-z}.
\]
To imply the excluded bound it needs uniform control of $r$ on
$L/\xi(g)\geq X$ for all sufficiently small $g$. Choose both
relative factors between $3/4$ and $5/4$ there. Their product is
within $[1/2,2]$. If $\xi(g)\leq Ag^{-\nu}$ for finite constants,
a larger integer power and rational coefficient give the required
computable threshold; an exponential envelope works in the same way.
A pointwise scaling ansatz without a uniform thermodynamic tail
does not imply this statement. The finite-size scaling hypotheses of
Fisher--Barber \cite{fisherBarber} are therefore not being disproved
in their ordinary applications. The result excludes extending this
specified uniform ansatz to these constructed gaps.

Size-driven transitions \cite{sizeDriven} already exhibit enormous
hidden scales in two-dimensional systems. Their constant-gap changes
of ground-state character across a family differ from the relative
gap-settling obstruction for one fixed family as coupling varies
here. Neither finite checks of particular couplings nor that prior
phenomenon replace the arithmetic proof of the present obstruction.
